\chapter{Síntese dos achados}

\bloco{O que as fichas sustentam}

\begin{enumerate}[leftmargin=*]
\item \textbf{O objeto do desacordo muda a cada marco.} Em 1906 a disputa é o
padrão, a quebra ou não da paridade de 1846, e o próprio mecanismo. Em 1910 é o
número da taxa e o tamanho do limite dos depósitos. Em 1914 é a emissão de
papel-moeda e o troco das notas. Serzedelo Corrêa, Barbosa Lima e Paula Ramos
se opõem à criação e depois defendem a letra da lei de 1906 contra a ampliação
do limite. Em várias trajetórias isso é mudança de objeto, não de lado (fichas
1 e 5).

\item \textbf{A oposição de 1906 não é um bloco doutrinário.} Convivem os que
querem o par de 1846 (Barbosa Lima, Paula Ramos), um defensor da quebra
definitiva em outra data e a outra taxa (Serzedelo Corrêa) e um estabilizador a
taxa mais baixa que a do projeto (Alcindo Guanabara). Serzedelo inclui a
valorização do café no próprio programa (ficha 1).

\item \textbf{Os defensores da Caixa não se descrevem como partidários do
câmbio baixo.} Campista quer a taxa corrente e a estabilidade como bem em si;
Carvalhal recusa o rótulo em 1910; Cincinato Braga quer a alta por degraus.
Codificar a defesa da Caixa como defesa do câmbio baixo seria erro de
construto (fichas 2 e 5).

\item \textbf{A Caixa foi criada pelo Congresso contra a posição declarada do governo que
terminava.} Campista o diz na tribuna, a imprensa o registra, Joaquim Murtinho
renuncia ao Senado por causa do projeto e Pinheiro Machado reconta o episódio
em 1913 (ficha 3).

\item \textbf{Alcindo Guanabara não é caso de reviravolta.} Vota a favor do
princípio em 1906, pede quebra a 12, defende 15 em 1910 e repete em 1913 que a
taxa certa era 12. A divergência aparente entre seus discursos e seu voto se
explica pelas emendas de 24 de setembro de 1906 (fichas 1, 5 e 6).

\item \textbf{Em 1910 o eixo é 15 contra a alta, com 16 como transação.} A
divisão acompanha a oposição que adoto como eixo do período:
valorização da moeda contra estabilidade a taxa nova (ficha 5).

\item \textbf{Os dois campos nomeiam interesses de classe.} Barbosa Lima em
1906, Bulhões, Murtinho e Calógeras em 1910 atribuem a defesa da taxa baixa a
lavradores, detentores de café ou a uma minoria; Cincinato Braga inverte os
termos e fala em defraudação das classes camponesas. São enquadramentos
explícitos, candidatos a código de interesses invocados (fichas 1 e 5).

\item \textbf{O argumento do depositário tem história anterior a 1914.}
Aparece na voz de \textit{O Paiz} em 1909, em Bulhões em 1910 como defesa do
direito de retirar o ouro, e em 1914 como condenação da suspensão do troco. Em
dezembro de 1914 ele convive, na fala de Serzedelo, com o desejo de ver a
Caixa acabar (fichas 4, 5 e 7).

\item \textbf{Há voz própria de jornal contra a Caixa em 1906.} \textit{O
Paiz}, de 24 de setembro a 10 de novembro, e o \textit{Correio da Manhã}, em
setembro e outubro, escrevem contra o projeto. O registro corrobora a memória
de Pinheiro Machado para esses dois jornais e não equivale a medida de posição
editorial (ficha 8).
\end{enumerate}

\bloco{O que continua em aberto}

\begin{itemize}[leftmargin=*,nosep]
\item O conteúdo da inimizade de Cincinato Braga à Caixa em 1906, que ele
declara em 1910 e que nenhum texto de 1906 documenta.
\item A posição de Rui Barbosa: as divergências que o \textit{Correio da Manhã}
lhe atribui em outubro de 1906 e a ausência de fala direta sobre a emissão em
1914.
\item O discurso de Bulhões em banquete e a mensagem de Rodrigues Alves contra
o Convênio de Taubaté em 1906, citados pelo \textit{Correio da Manhã} e não
localizados.
\item A autoria da série de \textit{O Paiz} contra a Caixa e a posição do
jornal entre março e agosto de 1906.
\item O espaço em branco no \textit{Correio da Manhã} de 28 de dezembro de
1913, no lugar em que \textit{O Paiz} imprime o próprio nome.
\item A natureza dos depósitos em Londres mencionados na reunião de agosto de
1913.
\end{itemize}

\bloco{Limites do material}

Os \textit{Documentos Parlamentares} cobrem apenas 1906 e 1910. Para os demais
anos as fichas dependem de imprensa, em boa parte resumos de sessão, que são
fala de terceiro. A busca na imprensa partiu das páginas que nomeiam a Caixa de
Conversão; discursos sobre câmbio e moeda que não a nomeiam ficaram fora. O
texto dos jornais foi lido no OCR da Biblioteca Nacional, cujo ruído varia por
jornal e por ano, e a camada corrigida dos volumes parlamentares ainda não tem
rotina reproduzível no repositório. As citações marcadas com adaga devem ser
conferidas na imagem antes de qualquer uso no texto da dissertação.
